\documentclass{exam}
\usepackage{../../shah311}

\hypersetup{colorlinks=true, linktoc=section, linkcolor=blue}

\pagestyle{headandfoot}
\firstpageheadrule
\runningheadrule
\firstpageheader{Prof. Rong \\ Real Analysis}{Homework 4}{Jeevan Shah}
\runningheader{Real Analysis \\ Homework 4}{}{Shah}
\firstpagefooter{}{}{}
\runningfooter{ }{\thepage}{ }

\printanswers
% hw4: pp59: 2.5.1, 2.5.2, 2.2.5, 2.5.9, 

% pp70: 2.6.2, 2.6.4, 2.7.2, 2.7.10
\begin{document}
\colorbox{red}{\underline{2.5.1}:}
\begin{solution}
    \begin{parts}
        \part This is impossible by the Bolzano-Weierstrass Theorem which states that every bounded sequence contains at least one convergent subsequence

        \part Consider the sequence 
        \[
            x_n = \left(\frac{1}{2}, \frac{1}{2}, \frac{1}{3}, \frac{2}{3}, \ldots, \frac{1}{k}, \frac{k-1}{k}\right).
        \]
        Then the subsequence $(x_{n_{k}}) = (x_{2n - 1})$ converges to $0$ and the subsequence $(x_{n_{k}}) = (x_{2n})$ converges to $1$ but neither $0$ or $1$ is in the original sequence $(x_n)$.

        \part Consider the sequence 
        \[
            x_n = \left(1, 1, \frac{1}{2}, 1, \frac{1}{2}, \frac{1}{3}, 1, \frac{1}{2}, \frac{1}{3}, \frac{1}{4}, \cdots \right).
        \]
        Then, every number in the form $\frac{1}{k}$ appears infinitely many times in $(x_n)$. To create a subsequence that converges to $\frac{1}{k}$ we simply choose the constant sequence $(x_{n_k}) = \frac{1}{k}$.

        \part This is impossible as any sequence that satisfies the requirements of containing subsequences that converge to every points in the infinite set $\braces{1, 1/2, 1/3, 1/4, 1/5, \ldots}$ will also contain a subsequence that converges to $0$.
    \end{parts}
\end{solution}

\colorbox{red}{\underline{2.5.2}:}
\begin{solution}
    \begin{parts}
        \part This is true for if every proper subsequence of $(x_n)$ converges then the subsequence that does not contain $x_1$ and the subsequence that does not contain $x_2$ must both converge. Since both of these sequences converge and contain infinitely many terms from the original series, they must converge to the same limit. It follows that $(x_n)$ must also converge to the same limit since limits are unique. 

        \part This is true. If $(x_n)$ contains a divergent subsequence, then the subsequence must be unbounded. It follows that $(x_n)$ is also unbounded and thus cannot be convergent since all convergent sequences are bounded. 

        \part This is true. If $(x_n)$ is bounded and still divergent, then for large $n$, $(x_n)$ must oscillate between two values. This implies that there is a subsequence converging to each.

        \part This is true. If $(x_n)$ contains a convergent subsequence, say $(x_{n_k}) \to L$ then $L$ must be a bound for the entire sequence $(x_n)$. It follows that since $(x_n)$ is bounded and monotone, it is also convergent.
    \end{parts}    
\end{solution}

\colorbox{red}{\underline{2.5.5}:}
\begin{solution}
    \begin{proof}
        If every subsequence of $(a_n)$ converges to some $a \in \RR$ then for every subsequence $(a_{n_{k_{i}}})$ there exists an $N_i \in \NN$ such that $|a_{n_{k_{i}}} - a| < \epsilon$ for all $k_i \geq N_i$. Now, choose $N = \max\braces{N_i \mid i \in \NN}$. It follows that $|a_n - a| < \epsilon$ for all $n \geq N$ and $\epsilon > 0$. Thus $(a_n) \to a$.
    \end{proof}    
\end{solution}

\colorbox{red}{\underline{2.5.9}:}
\begin{solution}
    \begin{proof}
        Since $s = \sup S$ it follows, by definition of supremum, that $s - \epsilon \in S$ and $s + \epsilon \not\in S$. So, by definition of $S$, an infinite amount of $a_n$ are greater than $s - \epsilon$ while only a finite amount of $a_n$ are greater than $s + \epsilon$. Thus, we can create the desired subsequence by choosing each $a_{n_k}$ such that $s - \frac{1}{k} < a_{n_{k}} < s + \frac{1}{k}$. Because there are infinitely many terms in the interval $\left(s - \frac{1}{i}, s + \frac{1}{i}\right]$ so we may always find $n_{k+1} > n_k$ for $a_{n_{k}} \in \left(s - \frac{1}{k}, s + \frac{1}{k}\right]$ and $a_{n_{k+1}} \in \left(s - \frac{1}{k+1}, s + \frac{1}{k+1}\right]$. Then, by the squeeze theorem, the limit of this subsequence must be $s$.
    \end{proof}    
\end{solution}

\colorbox{red}{\underline{2.6.2}:}
\begin{solution}
    \begin{parts}
        \part We may consider the alternating harmonic series, $a_n = \frac{(-1)^n}{n}$. This is a convergent series, thus Cauchy, but it is not montone.

        \part This is impossible since if a sequence $(a_n)$ is Cauchy, then it is also convergent and thus bounded. It is then impossible to have an unbounded subsequence within a bounded sequence.

        \part This is impossible since if a sequence  $(a_n)$ contains a Cauchy subsequence, $(a_n)$ must be bounded by the limit of the subsequence. However, if $(a_n)$ is divergent and monotone, it is then unbounded, which is a contradiction. 

        \part This is impossible since if a sequence $(a_n)$ contains a Cauchy subsequence, then $(a_n)$ must also be bounded by the limit of the subsequence. 
    \end{parts}
\end{solution}

\colorbox{red}{\underline{2.6.4}:}
\begin{solution}
    \begin{parts}
        \part If $(a_n)$ and $(b_n)$ are Cauchy, then we may say $(a_n) \to a$ and $(b_n) \to b$. It follows, by the Algebraic Limit Theorem, that $(c_n) \to |a - b|$. Since $(c_n)$ is a convergent sequence, it is Cauchy.

        \part Since $(a_n)$ is Cauchy we may say that $(a_n) \to a$. Then, for $n$ large, $c_n$ will oscillate between $-a$ and $a$. It follows that $c_n$ is divergent and thus not Cauchy.

        \part If $c_n = [[a_n]]$ then for $m, n \in \NN$ we have 
        \[
            |[[a_n]] - [[a_m]]| \geq 1
        \]
        for $[[a_n]] \neq [[a_m]]$. Then there exists $\epsilon < 1$ such that 
        \[
            |c_n - c_m| > \epsilon
        \]
        so $c_n$ cannot be Cauchy.
    \end{parts}    
\end{solution}

\colorbox{red}{\underline{2.7.2}:}
\begin{solution}
    \begin{parts}
        \part Since 
        \[
            \sum_{n=1}^\infty \frac{1}{2^n + n} \leq \sum_{n=1}^{\infty} \frac{1}{2^n}
        \]
        and $\sum \frac{1}{2^n}$ is a convergent series, the original series must also converge.

        \part Since $\sin(n) \leq 1$ for all $n \in \NN$ we have that 
        \[
            \sum_{n=1}^{\infty} \frac{\sin(n)}{n^2} \leq \sum_{n=1}^{\infty} \frac{1}{n^2} 
        \]
        but $\sum \frac{1}{n^2}$ is a convergent $p$-series. Thus, the original series must also converge.

        \part We can consider the closed form of the series 
        \[
            \sum_{n=1}^{\infty} \frac{(-1)^n(n+1)}{2n} 
        \]
        to see that this is an alternating series. Applying the alternating series test we can see that 
        \[
            \lim_{n\to\infty} a_n = \lim_{n\to\infty}\frac{n+1}{2n} = \frac{1}{2} \neq 0.
        \]
        Thus, the original series diverges.

        \part Consider the grouping of terms 
        \[
            \left(1 + \frac{1}{2} - \frac{1}{3}\right) + \left(\frac{1}{3} + \frac{1}{4} + \frac{1}{5}\right) + \cdots + \left(\frac{1}{3n - 2} + \frac{1}{3n - 1} - \frac{1}{3n}\right) + \cdots.
        \]
        Since 
        \[
            \frac{1}{3n-1} - \frac{1}{3n} > 0 
        \]
        because $3n - 1 < n$, it follows that 
        \[
            \frac{1}{3n-2} + \frac{1}{3n - 1} - \frac{1}{3n} > 0.
        \]
        Thus, every grouping is a positive number and so the sequence of partial sums must be unbounded. It follows that the series is divergent.

        \part Notice that for the odd numbered terms we have the series 
        \[
            \sum_{n=1}^{\infty} \frac{1}{2n-1} 
        \]
        while for the even term we have 
        \[
            \sum_{n=1}^{\infty} \frac{-1}{(2n)^2}.
        \]
        Since $\sum_{n=1}^{\infty} \frac{1}{2n-1}$ is divergent by being a variant of the harmonic series, the entire original must also be divergent.
    \end{parts}    
\end{solution}

\newpage 

\colorbox{red}{\underline{2.7.10}:}
\begin{solution}
    \begin{parts}
        \part Consider the closed form 
        \[
            \prod_{n=1}^{\infty} \frac{2^{n-1} + 1}{2^{n-1}} = \prod_{n=1}^{\infty} \left(1 + \frac{1}{2^{n-1}}\right).
        \]
        From exercise $(2.4.10)$ we know that product will converge if and only if 
        \[
            \sum \frac{1}{2^n-1} 
        \]
        converges. But, this is a geometric series with $|r| < 1$ and thus is convergent. So, the infinite product must also be convergent.

        \part The product converges since the sequence of partial products is bounded (by $0$ and $1$) and decreasing since we are always multiplying by a number less than $1$. Now, we consider the form 
        \[
            P = \prod_{n=1}^{\infty} \frac{2n-1}{2n} = \prod \left(1 - \frac{1}{2n}\right)
        \]
        and its reciprocal 
        \[
            P' = \frac{1}{P} = \prod_{n=1}^{\infty} \left(1 + \frac{1}{2n-1}\right).
        \]
        Using the results from exercise $(2.4.10)$ we know that $P'$ converges if, and only if, 
        \[
            \sum_{n=1}^{\infty} \frac{1}{2n-1} 
        \]
        converges. Since it diverges we have that $P' \to \infty$. But then 
        \[
            P = \frac{1}{P'} \Rightarrow P \to \frac{1}{\infty} = 0. 
        \]
        So the original product does converge to $0$.

        \part Consider the closed form 
        \[
            \prod_{n=1}^\infty \frac{(2n)^2}{(2n-1)(2n+1)} = \prod \frac{4n^2}{4n^2 - 1} = \prod \left(1 + \frac{1}{4n^2 - 1}\right).
        \]
        Just as before, since 
        \[
            \sum_{n=1}^{\infty} \frac{1}{4n^2 - 1} 
        \]
        is a convergent $p$-series, the infinite product must also be convergent.
    \end{parts}
\end{solution}
\end{document}